\minitoc

% =============================================================================
% Fondements Theoriques et Cadre de Risque

\section{Fondements Théoriques et Risques Systémiques}

L'apprentissage automatique profond repose sur des principes d'optimisation numérique empirique dont la robustesse n'est pas inconditionnelle.

\subsection{Le théorème du \emph{No Free Lunch}}

Formulé par David Wolpert (1996), le théorème du \emph{No Free Lunch} stipule qu'aucun algorithme d'apprentissage n'est universellement supérieur à tous les autres sur l'ensemble des distributions de données imaginables :
\begin{equation}
    \sum_{\mathcal{D}} \mathbb{E}\left[\mathcal{E}_{\mathrm{test}} \mid \mathcal{A}_1, \mathcal{D}\right] = \sum_{\mathcal{D}} \mathbb{E}\left[\mathcal{E}_{\mathrm{test}} \mid \mathcal{A}_2, \mathcal{D}\right].
\end{equation}
L'efficacité d'un modèle (ex: l'invariance par translation d'un CNN ou la dépendance temporelle ordonnée d'un RNN) repose exclusivement sur l'adéquation de ses \textbf{biais inductifs} avec les propriétés physiques du problème traité.

\subsection{Cadre légal européen : L'AI Act}

Face à l'opacité inhérente aux modèles dits "boîte noire", l'Union Européenne a introduit en 2024 un cadre législatif gradué selon quatre niveaux de criticité :
\begin{itemize}
    \item \textbf{Risque inacceptable} : interdiction absolue (ex : évaluation sociale généralisée, manipulation cognitive subliminale, identification biométrique en temps réel dans l'espace public à des fins policières).
    \item \textbf{Haut risque} : conformité, traçabilité et explicabilité strictes (ex : systèmes de tri automatisé de CV, justice prédictive, calcul des aides sociales).
    \item \textbf{Risque limité} : obligations de transparence (ex : systèmes génératifs d'images / deepfakes, agents conversationnels / chatbots).
    \item \textbf{Risque minimal} : usage libre (ex : filtres anti-spam, jeux vidéo).
\end{itemize}

% =============================================================================
% Biais et Equite Algorithmique

\section{Biais de Données et Métriques d'Équité (\emph{Fairness})}

\subsection{Typologie des biais dans le cycle de vie du modèle}

Un algorithme n'est jamais neutre : il optimise une fonction de coût sur des distributions historiques parfois biaisées.
\begin{enumerate}
    \item \textbf{Biais d'échantillonnage (Sampling bias)} : illustré historiquement par les travaux d'Abraham Wald sur les bombardiers de la Seconde Guerre mondiale (l'observation des impacts sur les avions revenus omettait les zones critiques des avions abattus).
    \item \textbf{Biais historique et stéréotypage} : propagation des corrélations fallacieuses présentes dans la langue écrite (ex: l'analogie vectorielle dans Word2Vec où $\vec{v}_{\text{docteur}} - \vec{v}_{\text{homme}} + \vec{v}_{\text{femme}} \approx \vec{v}_{\text{infirmière}}$).
    \item \textbf{Biais de mesure et disparités intersectionnelles} : travaux de Buolamwini et Gebru (2018) démontrant que les classifieurs commerciaux présentaient un taux d'erreur inférieur à $1\%$ pour les hommes à peau claire, mais supérieur à $34\%$ pour les femmes à peau sombre en raison de bases d'entraînement déséquilibrées.
\end{enumerate}

\subsection{Critères mathématiques d'équité}

Soit $X$ les attributs descriptifs, $A \in \{0, 1\}$ un attribut sensible protégé (genre, origine), $Y \in \{0, 1\}$ la vérité terrain et $\hat{Y} = f(X, A)$ la décision binaire du modèle :
\begin{definition}[Mesures formelles d'équité de groupe]
    \begin{itemize}
        \item \textbf{Parité statistique (Démographique)} : l'accès à la classe favorable est indépendant du groupe :
        \begin{equation}
            \mathbb{P}(\hat{Y} = 1 \mid A = 0) = \mathbb{P}(\hat{Y} = 1 \mid A = 1).
        \end{equation}
        \item \textbf{Égalité des chances (Equal Opportunity)} : le taux de vrais positifs est identique entre les groupes :
        \begin{equation}
            \mathbb{P}(\hat{Y} = 1 \mid A = 0, Y = 1) = \mathbb{P}(\hat{Y} = 1 \mid A = 1, Y = 1).
        \end{equation}
    \end{itemize}
\end{definition}

% =============================================================================
% Robustesse et Exemples Adversariaux

\section{Robustesse et Exemples Adversariaux}

\subsection{Vulnérabilité aux perturbations imperceptibles}

En grande dimension, la géométrie des réseaux de neurones engendre des discontinuités locales brutales : une micro-perturbation vectorielle $\delta$, non discernable pour l'œil humain, peut faire basculer la classification vers une cible aberrante avec une confiance proche de $100\%$.

\begin{definition}[Problème d'optimisation d'une attaque adversariale]
    Étant donné un classifieur $f(x) = \arg\max_c p(c \mid x)$ et une cible voulue $y_{\mathrm{adv}} \neq y$, la recherche d'une attaque minimale sous contrainte de norme $\ell_p$ s'écrit :
    \begin{equation}
        \min_{\delta} \|\delta\|_p \quad \text{tel que} \quad f(x + \delta) = y_{\mathrm{adv}}.
    \end{equation}
\end{definition}

\begin{proposition}[Attaque FGSM (\emph{Fast Gradient Sign Method})]
    Goodfellow et al. (2014) ont démontré qu'une attaque efficace sous contrainte $\ell_\infty$ ($\|\delta\|_\infty \le \epsilon$) s'obtient analytiquement par une passe linéaire inverse sur le gradient de la perte :
    \begin{equation}
        x_{\mathrm{adv}} = x + \epsilon \cdot \operatorname{sign}\left( \nabla_x \mathcal{L}(\theta, x, y) \right).
    \end{equation}
\end{proposition}

\subsection{Implémentation d'une attaque FGSM sous PyTorch}

\begin{lstlisting}[language=Python, caption={Génération d'un exemple adversarial FGSM sous PyTorch}]
import torch
import torch.nn as nn

def fgsm_attack(image: torch.Tensor, epsilon: float, data_grad: torch.Tensor) -> torch.Tensor:
    # Recupere le signe des gradients de la perte par rapport aux pixels d'entree
    sign_data_grad = data_grad.sign()
    # Ajout de la perturbation directionnelle bornee par epsilon
    perturbed_image = image + epsilon * sign_data_grad
    # Re-ecretage pour conserver des valeurs de pixels physiques valides [0, 1]
    perturbed_image = torch.clamp(perturbed_image, 0, 1)
    return perturbed_image

# Exemple d'execution
model.eval()
input_img.requires_grad = True

output = model(input_img)
loss = criterion(output, target_label)
model.zero_grad()
loss.backward()

# Perturbation de l'entree
grad_x = input_img.grad.data
perturbed_data = fgsm_attack(input_img, epsilon=0.05, data_grad=grad_x)
\end{lstlisting}

% =============================================================================
% Explicabilite et Interpretabilite (XAI)

\section{Explicabilité et Interprétabilité (XAI)}

\subsection{Cartes de saillance (\emph{Saliency Maps})}

Pour une entrée $x_0$ et une classe d'intérêt $c$, la carte de saillance de Simonyan et al. (2013) mesure l'influence linéaire locale au premier ordre de chaque coordonnée d'entrée sur le score $S_c$ avant normalisation Softmax :
\begin{equation}
    M_{\mathrm{saliency}}(x_0) = \left| \nabla_x S_c(x) \big|_{x = x_0} \right|.
\end{equation}
Cette méthode permet notamment de détecter si un classifieur s'appuie sur des artéfacts fallacieux d'arrière-plan (effet dit "Clever Hans") plutôt que sur les caractéristiques intrinsèques de l'objet.

\subsection{Perturbations locales : La méthode LIME}

Ribeiro et al. (2016) ont proposé l'approche LIME (\emph{Local Interpretable Model-agnostic Explanations}). Le modèle complexe $f$ est considéré comme une boîte noire complète. Autour d'un point d'intérêt $x$, on échantillonne un ensemble de points perturbés $z'$, pondérés par une mesure de proximité $\pi_x(z) = \exp(-D(x, z)^2 / \sigma^2)$, puis on ajuste un modèle linéaire simple et interprétable $g(z') = w_g^\top z'$ en minimisant :
\begin{equation}
    \arg\min_{g \in \mathcal{G}} \sum_{z, z'} \pi_x(z) \left( f(z) - g(z') \right)^2 + \Omega(g)
\end{equation}
où $\Omega(g)$ force la parcimonie sur les coefficients, révélant directement la contribution marginale locale de chaque composante d'entrée.